\section{VANISH Protocol Architecture}
\label{sec:protocol}

The VANISH protocol consists of two primary components: a portable binary container specification (\texttt{.vfile}) and an operating system lifecycle enforcement service.

\subsection{Binary Container Layout and Overhead Bounds}
A \texttt{.vfile} container begins with a fixed binary framing segment totaling exactly $140$\,bytes, followed by variable-length authenticated segments:
\begin{enumerate}
    \item \textbf{Fixed Header (32 bytes):} Encodes magic bytes (\texttt{VAN1}), protocol version ($1$), cryptographic capability flags, and big-endian segment lengths.
    \item \textbf{Public Key (32 bytes):} The Ed25519 public key of the issuing sender.
    \item \textbf{AEAD Nonce (12 bytes):} The 96-bit initialization vector for AES-256-GCM encryption (zero-padded if unencrypted).
    \item \textbf{Digital Signature (64 bytes):} Detached Ed25519 signature covering the unified cryptographic preimage.
    \item \textbf{Manifest Segment:} Canonical CBOR (RFC 8949) map encoding file metadata, MIME type, issuance time $t_{\text{iss}}$, expiration timestamp $t_{\text{exp}}$, policy action, and optional recipient key encapsulations. Manifests are bounded to a maximum length of $4096$\,bytes, with filename strings truncated to $\le 255$ UTF-8 bytes. In evaluated samples, this segment spans $180$--$289$\,bytes.
    \item \textbf{Payload Segment:} The raw payload or AES-256-GCM ciphertext with appended 16-byte authentication tag.
\end{enumerate}

Thus, total container framing overhead is strictly bounded:
\begin{equation}
    \Omega_{\text{overhead}} = 140 + |\text{Manifest}_{\text{CBOR}}| \quad (320 \le \Omega_{\text{overhead}} \le 4236 \text{ bytes})
\end{equation}
In standard office and multimedia documents, $\Omega_{\text{overhead}}$ averages $320$--$429$\,bytes.

\subsection{Cryptographic Integrity Binding and Key Distribution}
To ensure that an attacker cannot alter flags, replace the manifest, or extend the expiration date, the signature covers a unified cryptographic preimage:
\begin{equation}
    \mathcal{M} = \text{Header}_{0:26} \mathbin{\Vert} \text{PubKey} \mathbin{\Vert} \text{Nonce} \mathbin{\Vert} \text{Manifest}_{\text{CBOR}} \mathbin{\Vert} \mathcal{H}(\text{Payload})
\end{equation}
where $\mathcal{H}$ represents SHA-256. The sender computes the digital signature $\sigma$ as:
\begin{equation}
    \sigma = \text{Ed25519\_Sign}(\text{SK}_{\text{Sender}}, \mathcal{H}(\mathcal{M}))
\end{equation}
Because the signature binds the first 26 bytes of the header, the public key, the nonce, the canonical CBOR manifest, and the payload hash, any change to timestamps, policy flags, or file content will cause verification to fail.

\textbf{Key Distribution via HPKE:} In Tier 2 (Encrypted Mode), payload confidentiality across untrusted networks is achieved via hybrid envelope encryption (RFC 9180 HPKE)~\cite{rfc9180hpke}. The sender generates an ephemeral Content Encryption Key ($K_{\text{CEK}} \leftarrow \mathcal{K}_{\text{AES-256}}$) and encrypts the payload via AES-256-GCM. The sender then encapsulates $K_{\text{CEK}}$ using the recipient's public key ($\text{PK}_{\text{Recv}}$):
\begin{equation}
    C_{\text{key}} = \text{HPKE\_Seal}(\text{PK}_{\text{Recv}}, K_{\text{CEK}})
\end{equation}
The encapsulated key $C_{\text{key}}$ is placed inside the CBOR manifest under an authorized recipient map. Upon ingestion, the recipient's OS daemon un-encapsulates $K_{\text{CEK}}$ using the host private key securely held within the hardware keystore (ARM TrustZone or TPM 2.0) and imports it into the keystore tagged with the container's unique \texttt{file\_id}. Volatile memory buffers holding $K_{\text{CEK}}$ are immediately overwritten with zeros.

\subsection{Formal Security Reductions}
We formalize the security guarantees of the protocol under standard cryptographic assumptions:

\newtheorem{theorem}{Theorem}

\begin{theorem}[Manifest Unforgeability]
Assuming Ed25519 is Existentially Unforgeable under Chosen Message Attacks (EUF-CMA) and SHA-256 is collision resistant, no probabilistic polynomial-time (PPT) adversary $\mathcal{A}$ without the private key $\text{SK}_{\text{Sender}}$ can construct a forged container $\mathcal{C}'$ with an altered expiration timestamp $t'_{\text{exp}} > t_{\text{exp}}$ that is accepted by an authorized verifier pinned to $\text{PK}_{\text{Sender}}$, except with negligible probability $\text{negl}(\lambda)$.
\end{theorem}

\begin{IEEEproof}
Suppose $\mathcal{A}$ produces a modified container $\mathcal{C}'$ with manifest $\text{Manifest}'$ containing $t'_{\text{exp}} \ne t_{\text{exp}}$. By canonical CBOR serialization, $\text{Manifest}' \ne \text{Manifest}$, yielding preimage $\mathcal{M}' \ne \mathcal{M}$. By the collision resistance of SHA-256, $\mathcal{H}(\mathcal{M}') = \mathcal{H}(\mathcal{M})$ occurs with probability at most $\text{negl}(\lambda)$. Thus, for $\mathcal{C}'$ to be accepted, $\mathcal{A}$ must output a valid signature $\sigma'$ on a fresh message $\mathcal{H}(\mathcal{M}')$ under $\text{PK}_{\text{Sender}}$. Any PPT adversary that succeeds with non-negligible advantage can be directly reduced to an algorithm breaking the EUF-CMA security of Ed25519.
\end{IEEEproof}

\begin{theorem}[Conditional Post-Expiry Confidentiality]
In Tier 2, assuming: (i) $K_{\text{CEK}}$ was not exfiltrated prior to $t_{\text{exp}}$, (ii) ephemeral session keys are zeroized from host memory, and (iii) the hardware keystore irreversibly purges $K_{\text{CEK}}$ when $t \ge t_{\text{exp}}$, any PPT adversary $\mathcal{A}$ inspecting physical storage remnants has negligible advantage in recovering the plaintext payload under the IND-CCA2 security of AES-256-GCM.
\end{theorem}

\begin{IEEEproof}
Upon reaching $t \ge t_{\text{exp}}$, the host keystore purges $K_{\text{CEK}}$ via hardware-backed register zeroization. Under the assumption that volatile memory was scrubbed and $K_{\text{CEK}}$ was never exported, an adversary possessing post-expiry storage remnants (including uncollected flash blocks) observes only ciphertext $C$ and nonce $N$. Since $K_{\text{CEK}}$ is uniformly distributed in $\{0, 1\}^{256}$ and inaccessible, recovering any bit of plaintext reduces directly to distinguishing AES-256-GCM ciphertexts without the key, where $\mathbf{Adv}^{\text{IND-CCA2}}_{\text{AES-GCM}}(\mathcal{A}) \le \text{negl}(\lambda)$.
\end{IEEEproof}
